%% ma: L12-B06-0047
%% lop: 12
%% chuong: 2
%% bai: 6
%% dang: 12.6.18
%% loai: TN4
%% muc: TH
%% dap_an: "D"
%% nguon: "CĐ2.1.pdf giáo viên gửi 27/10/2026 (M07, Lê Minh Khôi – Nguyễn Anh Khôi, Chương II Chủ đề 1 Vectơ trong không gian), Câu 7"
%% kien_thuc: "Ba vectơ đồng phẳng (chương trình cũ, không có trong SGK KNTT 12 Bài 6)"
%% trang_thai: cach-ly
%% ly_do: "Vượt chương trình: khái niệm ba vectơ đồng phẳng không có trong SGK KNTT 12 Bài 6 (giống câu cách ly L12-B06-0040); phương án A, B, C chỉ là điều kiện đủ khi m, n cố định. Gợi ý: bỏ."
\begin{ex}
Với hai số thực $m$, $n$ khác $0$, điều kiện để ba vectơ $\vec a$, $\vec b$, $\vec c$ (trong đó $\vec a$, $\vec b$ không cùng phương) đồng phẳng là
\choice{$\vec c=m\cdot\vec a+n\cdot\vec b$}{$\vec c=m\cdot\vec a-n\cdot\vec b$}{$\vec c+m\cdot\vec a+n\cdot\vec b=\vec 0$}{\True A, B, C đều đúng}
\loigiai{Ba vectơ đồng phẳng ($\vec a$, $\vec b$ không cùng phương) khi $\vec c$ biểu diễn được theo $\vec a$, $\vec b$; cả ba dạng A, B, C đều là $\vec c=x\vec a+y\vec b$.}
\end{ex}
